% eq_axiom_2.tex — Axiom 2 (Topo-Kinematic Isomorphism)
% Single source of truth. \input this wherever Axiom 2 is formally restated.
% Per canonical Scheme A: manuscript/vol_1_foundations/chapters/01_fundamental_axioms.tex:54-59.
%
% PROVENANCE: pre-2026-04-27 this file was titled "Axiom 2 — Fine Structure" and
% documented the calibration constants alpha and V_yield. Per axiom homologation
% 2026-04-27, those constants are derived (alpha from Axiom 4's EMT operating
% point; V_yield from Axiom 4's saturation activation threshold); they have been
% moved to eq_calibration_constants.tex. This file now states the canonical
% primitive Axiom 2 from Vol 1 Ch 1.
%
\begin{resultbox}{Axiom 2 — Topo-Kinematic Isomorphism}
Charge $q$ is a discrete geometric dislocation in the substrate network. The Burgers vector of this dislocation IS the lattice pitch $\ell_{node}$, so the fundamental dimension of charge is identical to length ($[Q] \equiv [L]$). The macroscopic scaling is the Topological Conversion Constant:
\begin{equation}
    \xi_{topo} \equiv \frac{e}{\ell_{node}} \quad \text{[Coulombs / Meter]}
    \label{eq:axiom2_xi_topo}
\end{equation}
Numerical value $\xi_{topo} \approx 4.149 \times 10^{-7}$ C/m; see \texttt{eq\_calibration\_constants.tex}.

\medskip\noindent
\textbf{Charge quantization} follows directly: dislocation Burgers vectors must respect the K4 lattice (Axiom 1), so charge is quantized in integer multiples of $e$. \textbf{Charge sign} is the handedness of the dislocation in the substrate's chiral $I4_1 32$ structure --- particles and antiparticles are the two parity-related Burgers vectors permitted by the lattice geometry.

\medskip\noindent
\textbf{Fractional charges (Witten effect).} The $\mathbb{Z}_3$ symmetry of Borromean three-strand topology splits a unit dislocation into three linked sub-dislocations carrying $\pm 1/3\, e$ and $\pm 2/3\, e$. This is the substrate-native origin of quark fractional charges.

\medskip\noindent
\textbf{Topological solitons of this dislocation field are particles:} the electron is a real-space $0_1$ unknot carrying a $(2,3)$ phase-space torus-knot winding portrait on the Golden Torus (the Q-factor relation $Q_{TANK} = 1/\alpha$ is a \emph{calibration identity}, NOT a first-principles derivation of $\alpha$; Vol 1 Ch 8, honest-$\alpha$ scope); the proton is a $6^3_2$ Borromean linkage with $(2,5)$ phase-space winding; and the odd-$q$ $(2,q)$ family \emph{postdicts} the $S=0$ $N/\Delta$ single-Regge ladder (Vol 2 Ch 2). The trefoil lives in phase space; the soliton lives in real space.

\medskip\noindent
\textbf{Ladder honesty (postdiction scope).} The proton at $(2,5)$ is the one bare-topology emergence hit ($+0.74\%$, zero baryon input; $-0.002\%$ after the contained thermal residual). The rest of the ladder is a \textbf{postdiction}, not a forward chord: the matched $N/\Delta$ deviations run out to $+4.5\%$ (at $(2,13)$), $(2,21)$ has no PDG state within $5\%$, and the ensemble ``6/6'' hit-rate is null-dominated (random nearest-mass matching is expected to hit all 6 within $3\%$). Strange baryons and all mesons are off-ladder (open GAP). The surviving structural chord is the odd-$q$/even-$q$ link-exclusion ($(2,4)$ is not a knot) plus the curved-ladder FORM --- not the ensemble precision.

\medskip\noindent
\textbf{Interaction leg (2026-07-03).} The \emph{interaction} between two windings is now engine-derived: two $(2,3)$ windings repel when like-handed and attract when unlike-handed, a Coulomb sign structure mediated by the gapped rotation ($\omega$) sector (\texttt{clm-wcoul2}, consistency-class; the sign structure is robust, the screening scale is a host-knob artifact --- KB \texttt{vol4/claim-quality.md}). Axiom 2 thus carries both the charge \emph{dictionary} and a derived winding-pair \emph{interaction}; the formation-route (genesis of a winding from a free precursor) remains open.

\medskip\noindent
\textbf{Namespace caveat.} The Topological Conversion Constant $\xi_{topo}$ above is distinct from two other $\xi$ symbols in the framework: (a) Vol 3 Ch~\ref{ch:gravity_and_yield} (Gravity and Yield) defines a Machian $\xi \sim 10^{38}$ (cosmological-scale impedance boundary, derived $G = \hbar c / (7 \xi m_e^2)$), and (b) the substrate-scale Cosserat prefactors $\xi_{K1}, \xi_{K2}$ governing the Cosserat micropolar moduli at K4 nodes, with $\xi_{K2}/\xi_{K1} = 12$ K4-symmetry-forced (Vol 1 Ch~\ref{ch:macroscopic_moduli}, Macroscopic Moduli). See Glossary §1.5 for the namespace de-collision.
\end{resultbox}
